\documentclass[11pt,a4paper]{article}
\usepackage{curso}
\begin{document}
\portada{Lección}{2.6}{Ergodicidad}

\noindent\textbf{Unidad 2.}\quad\textbf{Objetivos:} promedios temporales frente a de ensamble; autocorrelación y
$\tau_{\text{int}}$; barras de error correctas. \textbf{Prerrequisitos:} 0.5, 2.4.

\section{Ergodicidad}
Un proceso estacionario es ergódico si $\bar x_T=\frac1T\int_0^Tx\,dt\to\E_\pi[x]$. En la práctica exige observar durante
mucho más que el tiempo de mezcla: un pozo doble con barrera de $12\,\kB T$ no cruza en la ventana observada y su promedio
temporal es engañoso (animación \evS).

\section{Autocorrelación y tamaño efectivo}
$\rho(\tau)=C(\tau)/C(0)$. Para muestras discretas, $\tau_{\text{int}}=1+2\sum_{k\ge1}\rho_k$ y $n_{\text{ef}}=n/\tau_{\text{int}}$.
Para un OU muestreado, $\rho_k=r^k$ con $r=e^{-a\Delta t}$ y $\tau_{\text{int}}=\frac{1+r}{1-r}$.
\begin{resultado}
$\Var(\bar x)\approx\Var(x)\,\tau_{\text{int}}/n$, es decir $\sigma_{\bar x}\approx\sqrt{\Var(x)\tau_{\text{int}}/n}$.
\end{resultado}
Repitiendo 200 veces una serie AR(1) con $r=0{,}98$ ($\tau_{\text{int}}=99$), las barras ingenuas $2\sigma/\sqrt n$ cubren el valor
verdadero en una fracción muy inferior al 95\,\%, mientras que las corregidas cubren $\approx95\,\%$ \evS.
\textbf{Método de bloques:} dividir en bloques de longitud $\gg\tau_{\text{int}}$, promediar cada bloque y usar la dispersión
de las medias de bloque. Para estimar $\tau_{\text{int}}$, truncar la suma con una ventana (la cola es ruidosa).

\begin{ejemplo}
OU con $a=0{,}1$ leído cada $\Delta t=0{,}1$: $r=e^{-0{,}01}$, $\tau_{\text{int}}\approx200$; con $n=10^4$, $n_{\text{ef}}\approx50$ y el
error real es $\sqrt{200}\approx14$ veces el ingenuo.
\end{ejemplo}

\begin{aplicacion}
Tiempo de integración del solucionador OU ($\tau_{\text{int}}\sim1/\lambda_{\min}$, Unidad 6); leer pbits más rápido que su
correlación desperdicia E/S (Unidad 5); Gibbs cerca de $T_c$ (4.4); barras de error en todos los laboratorios.
\end{aplicacion}

\begin{error}
\begin{itemize}[nosep]
\item $\sigma/\sqrt n$ con muestras correlacionadas. \quad Olvidar el burn-in.
\item Ergodicidad teórica $\neq$ práctica. \quad Sumar toda la cola ruidosa de $\rho_k$.
\end{itemize}
\end{error}

\begin{thebibliography}{9}\small
\bibitem[13]{katz} H. G. Katzgraber, Introduction to Monte Carlo methods, 2009.
\bibitem[3]{levin} D. A. Levin, Y. Peres y E. L. Wilmer, \emph{Markov Chains and Mixing Times}, 2009.
\bibitem[30]{tla} M. Aifer et al., Thermodynamic linear algebra, 2024.
\end{thebibliography}
\end{document}
